\documentclass[capchap,sumarioincompleto,a4paper,12pt,ruledheader,normaltoc]{./sty/utfprtex}%

\usepackage{./sty/utfprtexEEPropostaTCC}
\usepackage[utf8]{inputenc}
\usepackage[pdftex,linktocpage=true,pdfauthor={Seu Nome},pdftitle={Teste de coisa},pdfsubject={Exemplo de documento},pdfcreator={LaTeX e Dvips}]{hyperref} 
\usepackage{hyphenat}
\tolerance=1
\emergencystretch=\maxdimen
\hyphenpenalty=10000
\hbadness=10000
\hyphenchar\font=-1
\sloppy

\usepackage[alf,abnt-emphasize=bf,abnt-etal-text=it,bibjustif,abnt-etal-cite=2,abnt-etal-list=0,abnt-full-initials=yes,abnt-repeated-author-omit=yes,abnt-url-package=hyperref,abnt-last-names=abnt,abnt-missing-year=void]{abntex2cite} %

% Pacote para inserção de figuras em formato .eps.
\usepackage{graphicx}

% Pacote com mais opções para equações
\usepackage{amsmath}

%\renewcommand{\familydefault}{cmr}

\pdfstringdefDisableCommands{\edef\uppercase{}}
\renewcommand{\floatpagefraction}{0.90}
\renewcommand{\topfraction}{0.95}
\renewcommand{\bottomfraction}{0.95}
\renewcommand{\textfraction}{0.05}
\setlength{\intextsep}{5pt}
\setlength{\textfloatsep}{5pt}
\setlength{\floatsep}{5pt}


%%%%%%%%% Início das definições iniciais %%%%%%%%%


\autorA{Nome do Autor}
\autorAdados{Endereço A\\ emailA@email.com\\ telefone A}


\titulo{Título} 
\local{PATO BRANCO}
\data{ \the\year }
%%%%%%%%% Fim das definições iniciais %%%%%%%%%

\begin{document} % Início do documento.

% Monta a capa do documento.
\capa

%==================================================================

\input{content}

% a próxima linha é um conjunto de exemplos. Vc pode comentá-la somente.
\input{exemplos}

%==================================================================
% REFERÊNCIAS BIBLIOGRÁFICAS
% Bibliografia do arquivo bibliografia.bib
\bibliographystyle{./sty/abnt-num} % Estilo autor-data
\bibliography{bibliografia}    % As referencias deste testo estão no arquivo bibliografia.bib
%\ABNTaddcontentsline{toc}{chapter}{Referências bibliográficas}

\end{document}

